\section{Tutorial 7: Building an MLP and Backpropagation}

\subsection{Objectives}

This tutorial introduces neural networks as concrete computational systems:

\begin{itemize}
    \item Construct a simple multi-layer perceptron (MLP)
    \item Perform forward computation step by step
    \item Derive gradients manually for a small network
    \item Understand the need for backpropagation
\end{itemize}

\medskip

\noindent
\textbf{Focus.}  
Understand how neural networks compute and how gradients flow through them.

\subsection{Exercise 1: A Simple Neural Network}

Consider a two-layer neural network:

\[
h = \sigma(W_1 x + b_1), \quad
\hat{y} = W_2 h + b_2.
\]

\medskip

\noindent
\textbf{Tasks.}
\begin{enumerate}
    \item Specify the dimensions of all variables.
    \item Describe the forward pass.
\end{enumerate}

\medskip

\noindent
\textbf{Solution.}

Let:
\begin{itemize}
    \item $x \in \mathbb{R}^d$
    \item $W_1 \in \mathbb{R}^{m \times d}$, $b_1 \in \mathbb{R}^m$
    \item $h \in \mathbb{R}^m$
    \item $W_2 \in \mathbb{R}^{1 \times m}$, $b_2 \in \mathbb{R}$
\end{itemize}

\medskip

\noindent
Forward pass:
\begin{itemize}
    \item Linear transformation: $z_1 = W_1 x + b_1$
    \item Nonlinearity: $h = \sigma(z_1)$
    \item Output: $\hat{y} = W_2 h + b_2$
\end{itemize}

\medskip

\noindent
\textbf{Key takeaway.}  
Neural networks alternate between linear transformations and nonlinearities.

\subsection{Exercise 2: Forward Computation (Concrete Example)}

Let:
\[
x = \begin{pmatrix}1 \\ 2\end{pmatrix}, \quad
W_1 =
\begin{pmatrix}
1 & -1 \\
0 & 2
\end{pmatrix}, \quad
b_1 = \begin{pmatrix}0 \\ 1\end{pmatrix}.
\]

Use ReLU activation: $\sigma(z) = \max(0,z)$.

\medskip

\noindent
\textbf{Tasks.}
\begin{enumerate}
    \item Compute $z_1$
    \item Compute $h$
\end{enumerate}

\medskip

\noindent
\textbf{Solution.}

\[
z_1 =
\begin{pmatrix}
1\cdot1 + (-1)\cdot2 \\
0\cdot1 + 2\cdot2
\end{pmatrix}
+
\begin{pmatrix}0 \\ 1\end{pmatrix}
=
\begin{pmatrix}-1 \\ 5\end{pmatrix}.
\]

\[
h = \begin{pmatrix}0 \\ 5\end{pmatrix}.
\]

\medskip

\noindent
\textbf{Insight.}  
Nonlinearity introduces piecewise behavior.

\subsection{Exercise 3: Loss Function}

Let the target be $y$ and define squared loss:
\[
L = \frac{1}{2} (\hat{y} - y)^2.
\]

\medskip

\noindent
\textbf{Task.}
\begin{enumerate}
    \item Express $L$ in terms of network parameters.
\end{enumerate}

\medskip

\noindent
\textbf{Solution.}

\[
L = \frac{1}{2} \left(W_2 \sigma(W_1 x + b_1) + b_2 - y\right)^2.
\]

\medskip

\noindent
\textbf{Insight.}  
The loss is a nested composition of functions.

\subsection{Exercise 4: Gradient with Respect to Output Layer}

\textbf{Tasks.}
\begin{enumerate}
    \item Compute $\frac{\partial L}{\partial W_2}$
    \item Compute $\frac{\partial L}{\partial b_2}$
\end{enumerate}

\medskip

\noindent
\textbf{Solution.}

Let $\hat{y} = W_2 h + b_2$.

\[
\frac{\partial L}{\partial \hat{y}} = \hat{y} - y.
\]

\[
\frac{\partial L}{\partial W_2} = (\hat{y} - y) h^\top.
\]

\[
\frac{\partial L}{\partial b_2} = \hat{y} - y.
\]

\medskip

\noindent
\textbf{Insight.}  
Gradients at the output layer are straightforward.

\subsection{Exercise 5: Backpropagation to Hidden Layer}

\textbf{Tasks.}
\begin{enumerate}
    \item Compute $\frac{\partial L}{\partial h}$
    \item Compute $\frac{\partial L}{\partial z_1}$
\end{enumerate}

\medskip

\noindent
\textbf{Solution.}

\[
\frac{\partial L}{\partial h} = (\hat{y} - y) W_2^\top.
\]

\medskip

\noindent
For ReLU:
\[
\frac{\partial h}{\partial z_1} =
\begin{cases}
1 & z_1 > 0 \\
0 & z_1 \leq 0
\end{cases}
\]

\[
\frac{\partial L}{\partial z_1} =
\frac{\partial L}{\partial h} \odot \sigma'(z_1).
\]

\medskip

\noindent
\textbf{Insight.}  
Gradients are modulated by activation derivatives.

\subsection{Exercise 6: Gradients for First Layer}

\textbf{Tasks.}
\begin{enumerate}
    \item Compute $\frac{\partial L}{\partial W_1}$
    \item Compute $\frac{\partial L}{\partial b_1}$
\end{enumerate}

\medskip

\noindent
\textbf{Solution.}

\[
\frac{\partial L}{\partial W_1} =
\frac{\partial L}{\partial z_1} x^\top.
\]

\[
\frac{\partial L}{\partial b_1} =
\frac{\partial L}{\partial z_1}.
\]

\medskip

\noindent
\textbf{Insight.}  
Gradients propagate backward through layers via the chain rule.

\subsection{Exercise 7: Computational Perspective}

\textbf{Question.}  
Why is naive differentiation inefficient for deep networks?

\medskip

\noindent
\textbf{Discussion.}

\begin{itemize}
    \item Repeated computation of intermediate derivatives
    \item Exponential growth in computation if not reused
\end{itemize}

\medskip

\noindent
\textbf{Key idea.}  
Backpropagation reuses intermediate results efficiently.

\subsection{Exercise 8 (Discussion): What Backpropagation Does}

\textbf{Question.}  
What is the core idea of backpropagation?

\medskip

\noindent
\textbf{Discussion points.}

\begin{itemize}
    \item Apply chain rule systematically
    \item Propagate gradients from output to input
    \item Reuse computations
\end{itemize}

\medskip

\noindent
\textbf{Insight.}  
Backpropagation is dynamic programming for gradient computation.

\subsection{Summary}

\begin{itemize}
    \item Neural networks compute via layered transformations
    \item Gradients are computed using the chain rule
    \item Backpropagation efficiently propagates gradients
    \item Training neural networks requires structured gradient computation
\end{itemize}

\medskip

\noindent
\textbf{Preparation for next lecture.}  
We now formalize backpropagation and automatic differentiation, enabling efficient training of deep networks.